\section*{Bab \ref{ch:g8:pythagoras} --- Teorema Pythagoras}

\begin{solution}{exo:g8:pythagoras:1}
$\sqrt{81 + 144} = \sqrt{225} = 15$; \quad
$\sqrt{25 + 144} = \sqrt{169} = 13$; \quad
$\sqrt{64 + 225} = \sqrt{289} = 17$.
\end{solution}

\begin{solution}{exo:g8:pythagoras:2}
Sisi siku-siku$^2 = 25^2 - 7^2 = 625 - 49 = 576$, jadi sisi
siku-sikunya adalah $\sqrt{576} = 24$.
\end{solution}

\begin{solution}{exo:g8:pythagoras:3}
Persegi panjang: $\sqrt{12^2 + 9^2} = \sqrt{144 + 81} = \sqrt{225} =
15$ cm.

Persegi: $d = 5\sqrt2 \approx 7.1$ cm
(\cref{ex:g8:pythagoras:diagonal}).
\end{solution}

\begin{solution}{exo:g8:pythagoras:4}
$(20, 21, 29)$: $29^2 = 841$ dan $20^2 + 21^2 = 400 + 441 = 841$:
siku-siku.

$(7, 8, 11)$: $121 \neq 49 + 64 = 113$: bukan siku-siku.

$(1.5, 2, 2.5)$: $6.25 = 2.25 + 4$: siku-siku (ini
$(3,4,5)$ yang dibagi dua).
\end{solution}

\begin{solution}{exo:g8:pythagoras:5}
Papannya $= \sqrt{3^2 + 1.6^2} = \sqrt{9 + 2.56} = \sqrt{11.56} = 3.4$
m.
\end{solution}

\begin{solution}{exo:g8:pythagoras:6}
Tingginya jatuh di titik tengah alasnya (karena \omterm{def:g6:shapes:triangles}{segitiga sama kaki}),
jadi tiap setengah segitiganya punya sisi miring $10$ dan satu sisi
siku-siku $6$: $h = \sqrt{100 - 36} = \sqrt{64} = 8$ cm. Luasnya:
$\frac{12 \times 8}{2} = 48$ cm$^2$.
\end{solution}

\begin{solution}{exo:g8:pythagoras:7}
Diagonalnya: $\sqrt{88.5^2 + 49.8^2} = \sqrt{7832.25 + 2480.04}
= \sqrt{10312.29} \approx 101.5$ cm. Dalam inci:
$101.5 \div 2.54 \approx 40$ inci.
\end{solution}

\begin{solution}{exo:g8:pythagoras:8}
Kedua sisi siku-sikunya $24$ km dan $10$ km:
$\sqrt{24^2 + 10^2} = \sqrt{576 + 100} = \sqrt{676} = 26$ km.
\end{solution}

\begin{solution}{exo:g8:pythagoras:9}
Jarak kakinya: $\sqrt{2.5^2 - 2.4^2} = \sqrt{6.25 - 5.76} =
\sqrt{0.49} = 0.7$ m. \omterm{def:g3:division:half}{Seperempat} panjang tangganya adalah
$2.5 \div 4 \approx 0.6$ m: $0.7$ m yang ditemukan dekat dengan letak
yang dianjurkan --- dapat diterima.
\end{solution}

\begin{solution}{exo:g8:pythagoras:10}
Sisi siku-sikunya berukuran $5 - 1 = 4$ (mendatar) dan $5 - 2 = 3$
(tegak), jadi
\[
AB = \sqrt{4^2 + 3^2} = \sqrt{25} = 5 .
\]
\end{solution}

\begin{solution}{exo:g8:pythagoras:11}
\emph{1.} Secara langsung: $(a + b)^2$. Sebagai kepingan: empat
segitiga yang luasnya $\frac{ab}{2}$ masing-masing, ditambah persegi
miringnya $c^2$:
\[
(a + b)^2 = 4 \times \frac{ab}{2} + c^2 = 2ab + c^2 .
\]

\emph{2.} Dengan menjabarkan ruas kirinya:
$(a+b)^2 = a^2 + 2ab + b^2$. Jadi
\[
a^2 + 2ab + b^2 = 2ab + c^2
\quad\Longrightarrow\quad
a^2 + b^2 = c^2 . \qed
\]
\end{solution}

\begin{solution}{pb:g8:pythagoras:1}
\textbf{1.} $5^2 + 12^2 = 25 + 144 = 169 = 13^2$; \quad
$8^2 + 15^2 = 64 + 225 = 289 = 17^2$; \quad
$20^2 + 21^2 = 400 + 441 = 841 = 29^2$. Ketiganya tripel, dan menurut
\cref{thm:g8:pythagoras:converse} segitiga yang panjang sisinya
demikian adalah siku-siku, dengan \omterm{def:g3:shapes:rightangle}{sudut siku-sikunya} menghadap sisi
terpanjangnya.

\textbf{2.} Talinya membentuk segitiga bersisi $3$, $4$, $5$
(dalam ruas), dan $3^2 + 4^2 = 25 = 5^2$: menurut konversnya
(\cref{thm:g8:pythagoras:converse}), sudut antara sisi
$3$ dan $4$ ruas itu tepat siku-siku --- tanpa busur derajat,
hanya dengan simpul.

\textbf{3.} Kalau $a^2 + b^2 = c^2$, maka
\[
(ka)^2 + (kb)^2 = k^2 a^2 + k^2 b^2
= k^2\left(a^2 + b^2\right) = k^2 c^2 = (kc)^2 .
\]
Dari $(3, 4, 5)$: misalnya $(6, 8, 10)$, $(9, 12, 15)$ dan
$(30, 40, 50)$.

\textbf{4.} \omterm{def:g5:division:multiple}{Kelipatan} $(3k, 4k, 5k)$ yang memuat $5$ sebagai bilangan
terkecilnya memerlukan $3k = 5$, dan tidak ada bilangan bulat $k$
yang memenuhinya ($k = 1$ memberi $3$, $k = 2$ memberi $6$). Jadi
$(5, 12, 13)$ bukan \omterm{def:g5:division:multiple}{kelipatan} $(3, 4, 5)$.

\textbf{5.} Pertanyaan 3 menunjukkan bahwa penskalaan hanya
menghasilkan tripel $(3k, 4k, 5k)$; pertanyaan 4 memperlihatkan tripel
sungguhan yang tidak berbentuk demikian. Jadi keluarga semua \omterm{pb:g8:pythagoras:1}{tripel
Pythagoras} benar-benar lebih besar daripada \omterm{def:g5:division:multiple}{kelipatan} $(3, 4, 5)$:
menghasilkan semuanya menuntut sesuatu yang baru.

\textbf{6.} Tulislah $A = m^2$ dan $B = n^2$. Menurut kesamaan
istimewa (\cref{pb:g8:equations:1}):
\[
a^2 = (A - B)^2 = A^2 - 2AB + B^2,
\qquad
b^2 = (2mn)^2 = 4m^2n^2 = 4AB .
\]
Dengan menjumlahkannya:
\[
a^2 + b^2 = A^2 - 2AB + B^2 + 4AB
= A^2 + 2AB + B^2 = (A + B)^2
= \left(m^2 + n^2\right)^2 = c^2 .
\]

\textbf{7.}
\begin{center}
\renewcommand{\arraystretch}{1.2}
\begin{tabular}{c|ccc}
$(m, n)$ & $m^2 - n^2$ & $2mn$ & $m^2 + n^2$ \\
\hline
$(2, 1)$ & $3$ & $4$ & $5$ \\
$(3, 1)$ & $8$ & $6$ & $10$ \\
$(3, 2)$ & $5$ & $12$ & $13$ \\
$(4, 1)$ & $15$ & $8$ & $17$ \\
$(4, 3)$ & $7$ & $24$ & $25$ \\
\end{tabular}
\end{center}

\textbf{8.} $(8, 6, 10)$ adalah dua kali $(4, 3, 5)$ --- yaitu dua
kali $(3, 4, 5)$ milik tukang bangunan. Mesinnya kadang menemukan
kembali tripel lama dalam samaran.

\textbf{9.} Kita perlu $m^2 - n^2 = 21$ dan $2mn = 20$, jadi
$mn = 10$: mencoba $m = 5$, $n = 2$ memberi $25 - 4 = 21$ dan
$2 \times 10 = 20$, dengan sisi miring $25 + 4 = 29$. Mesinnya
mengeluarkan $(21, 20, 29)$ untuk $(m, n) = (5, 2)$.

\textbf{10.} $m = 8$, $n = 6$ memberi $m^2 + n^2 = 64 + 36 = 100$.
Tripelnya adalah
\[
\left(64 - 36,\ 2 \times 48,\ 100\right) = (28,\ 96,\ 100),
\]
dan memang $28^2 + 96^2 = 784 + 9216 = 10\,000 = 100^2$.

\textbf{11.} Dengan $m = n + 1$:
\[
m^2 - n^2 = (n + 1)^2 - n^2 = 2n + 1
\]
(yaitu \omterm{ex:g1:subtraction:difference}{selisih} kuadrat yang berurutan, \cref{pb:g8:equations:1});
$2mn = 2n(n + 1) = 2n^2 + 2n$; dan
$m^2 + n^2 = n^2 + n^2 + 2n + 1 = 2n^2 + 2n + 1$. Sisi miringnya
melampaui sisi siku-siku yang \omterm{def:g3:numbers:evenodd}{genap} tepat $1$. Untuk $n = 1, 2, 3, 4$:
\[
(3, 4, 5), \quad (5, 12, 13), \quad (7, 24, 25), \quad
(9, 40, 41) .
\]

\textbf{12.} Setiap \omterm{def:g3:numbers:evenodd}{bilangan ganjil} yang sedikitnya $3$ berbentuk
$2n + 1$ untuk suatu $n \geq 1$, dan pertanyaan 11 menyediakan tripel
yang sisi siku-siku \omterm{def:g3:numbers:evenodd}{ganjilnya} persis $2n + 1$. Untuk
$11 = 2 \times 5 + 1$, ambillah $n = 5$:
\[
(11,\ 60,\ 61),
\qquad
11^2 + 60^2 = 121 + 3600 = 3721 = 61^2 .
\]

\textbf{13.} Dengan $n = 1$ mesinnya memberi
$\left(m^2 - 1,\ 2m,\ m^2 + 1\right)$, yang sisi siku-siku \omterm{def:g3:numbers:evenodd}{genapnya}
$2m$ dapat menjadi \omterm{def:g3:numbers:evenodd}{bilangan genap} apa pun mulai dari $4$. Untuk
$14 = 2 \times 7$, ambillah $m = 7$: tripelnya $(48, 14, 50)$. Tripel
yang sama muncul dengan menggandakan $(7, 24, 25)$ dari tabel
pertanyaan 7: $(14, 48, 50)$ --- ketiga bilangan yang sama.

\textbf{14.} \omterm{def:g3:numbers:evenodd}{Bilangan genap} berbentuk $2k$, dan $(2k)^2 = 4k^2$ itu
\omterm{def:g3:numbers:evenodd}{genap} (\omterm{def:g5:division:multiple}{kelipatan} $4$, jadi \omterm{def:g3:numbers:evenodd}{genap}). \omterm{def:g3:numbers:evenodd}{Bilangan ganjil} berbentuk
$2k + 1$, dan kesamaan istimewa yang pertama memberi
\[
(2k + 1)^2 = 4k^2 + 4k + 1 = 2\left(2k^2 + 2k\right) + 1 ,
\]
yaitu \omterm{def:g3:numbers:evenodd}{bilangan genap} ditambah $1$: jadi \omterm{def:g3:numbers:evenodd}{ganjil}.

\textbf{15.} Andaikan $a$, $b$, $c$ semuanya \omterm{def:g3:numbers:evenodd}{ganjil}. Menurut
pertanyaan 14, $a^2$ dan $b^2$ \omterm{def:g3:numbers:evenodd}{ganjil}, jadi $a^2 + b^2$ adalah \omterm{def:g1:addition:def}{jumlah}
dua \omterm{def:g3:numbers:evenodd}{bilangan ganjil}: \emph{\omterm{def:g3:numbers:evenodd}{genap}}. Tetapi $c^2$ akan \emph{\omterm{def:g3:numbers:evenodd}{ganjil}}.
Kesamaan $a^2 + b^2 = c^2$ kalau begitu mustahil. Jadi pada setiap
\omterm{pb:g8:pythagoras:1}{tripel Pythagoras} sedikitnya satu dari ketiga bilangannya \omterm{def:g3:numbers:evenodd}{genap} ---
seperti yang dibenarkan setiap baris tabel pertanyaan 7.

\end{solution}
